\documentclass[11pt]{article}
\usepackage{amsmath,hyperref}
\title{Two passage mouths and a distinct contact boundary in Cave Circuit 01}
\author{}\date{}
\begin{document}\maketitle

\paragraph{Scope.}
The inherited development mission treats an edge of the published
Cave Circuit 01 visibility graph as a unit movement. That is a
deliberate abstraction, not a certificate of travel or a simulation of
a rock striking a passage. This note evaluates whether the chosen edge
has a source-derived mouth and whether another independently specified
edge removal can eliminate the return route. The world is already
viewed; the independent evaluation geometry and event order stay unspent.

\paragraph{From tiles to mouth coordinates.}
The pinned \href{https://github.com/osrf/subt/blob/083c6a8acb0c25aa1a35e34517e628cd3b52ae77/subt_ign/src/ConnectionHelper.cc}
{connector helper} transforms each tile's type-local connection points
by the tile pose, tests three coordinate differences against the
componentwise tolerance $1.5$, and uses the first match.
The pinned \href{https://github.com/osrf/subt/blob/083c6a8acb0c25aa1a35e34517e628cd3b52ae77/subt_ign/src/apps/dot_generator.cc}
{generator} assigns weights from the tile types, not from estimated
seconds of travel. The replay recomputes 56,750 port-pair tests among the
158 labeled staging and tile vertices; it matches all 163 published
weighted edges with one pair of mouths per edge. The largest within-pair
coordinate discrepancy is below $10^{-4}$ m. The twelve named caps
are not graph vertices and are separately checked not to match any
labeled port at that tolerance. No collision asset geometry or grid lookup
is expanded in this computation.

For the short return edge $2\!-\!3$, both source mouth positions are
near $(62.5,50,0)$ m. For the alternate route's $13\!-\!42$ edge,
they are near $(100,62.5,0)$ m. The fourth rockfall detector box is
centered at $(99.3555,62.4272,7.54495)$ with side lengths $(16,25,15)$
m. The alternate mouth is only $0.04495$ m below its bottom;
the short mouth is over $28.85$ m from the box. At the alternate
mouth center, either endpoint of the matched pair intersects the
detector for an axis-aligned performer box whose vertical half-extent
is at least $0.04495$ m, and not for a smaller box centered there.
This is a statement about the declared box with that sampled center,
not a requirement that a moving robot sample the center. The pinned
\href{https://github.com/gazebosim/gz-sim/blob/2a076e91c6df7699284ab93e244894bb0f55336f/src/systems/performer_detector/PerformerDetector.cc}
{detector} evaluates the performer component box at the parent
position, and its fourth SDF publisher matches entry. Intersection
is not evidence of a delivered deployment command, rock settling,
or blocked traversal. The other three detector boxes' matched mouth
locations are listed by the replay.

\paragraph{Two distinct failure consumers.}
With the short edge removed, a courier can still reach dock 39 from
scout 3 via $8,6,16,42,13,1,43$. Removing only $13\!-\!42$ leaves the
short path through tile 2. Removing both traps tile 3 in a 78-vertex
component of the imported movement graph, with neither dock 39 nor
staging 0 reachable. A second walk directly from the published labels
independently agrees on all four reachability verdicts. The two-edge
cut is minimal by inclusion; these are not the only candidate cuts.
Report custody and vehicle egress must then be handled separately:
an authorized holder across the cut can deliver an already surviving
copy, while that cannot extract the stranded vehicle. No refuge is
presumed. A named fallback still needs modeled and tested reachability.

The \href{https://github.com/osrf/subt/blob/083c6a8acb0c25aa1a35e34517e628cd3b52ae77/subt_ign/src/VisibilityRfModel.cc}
{pinned visibility RF computation} uses the sender/receiver poses,
the loaded graph and coordinate lookup, representative breadcrumb
positions and RF configuration. It does not use the detector
predicate or a newly placed collision body as a ray blocker. In the
loaded nominal graph without breadcrumbs the weighted visibility cost
between tile labels 3 and 39 is 12, assuming the queried coordinates
are classified into those labels. If one incorrectly used the edited
motion graph for RF as well, the cost would be 25 after the short
closure and the pair would be disconnected after both removals.
These are weighted graph costs, not measured packet outcomes.
Because the fourth event is triggered by approach, it can affect
physical passage and robot motion, which may in turn affect radio.
It cannot be counted as direct radio occlusion in this evaluator.
For a matched joint-fault test, establish contact loss separately
(e.g.\ a separately instantiated relay/pose change with the same
loaded LUT), and account for any common cause explicitly.

\paragraph{Next experiment and falsifier.}
Record the local approach clearance, actual motion, named report key,
authority and settlement, effective RF inputs and contact in a
single history. The first useful development comparison checks route
closure only, contact loss only, their combination, and a second
closure or reserve failure before the named receipt settles. Compare
a timed return-and-rendezvous or data-mule baseline that may use the
same observations, authority and live reserve; the unmodified
weighted graph is not a fair route oracle. Require a positive
admitted history with inspected work, a delivered report, and
egress or established fallback; the two-edge witness is a
boundary control, not a positive. If a vehicle can still bypass a
supposedly blocked mouth, or the observed pose is never classified
into the asserted label, reject that source bridge rather than
counting the nominal cut as a deployed effect. No independent
mission acceptance or operational superiority is shown here.
\end{document}
